The kernel consumes Bernoulli rates. Estimating those rates from data is a
separate statistical problem. For one Bernoulli rate estimated from $n$ samples,
Hoeffding's inequality gives
\[
  \Prob(|\hat p-p| > \epsilon) \le 2\exp(-2n\epsilon^2).
\]
Thus a two-sided radius at error probability $\delta$ is
\[
  \epsilon(n,\delta) = \sqrt{\frac{\log(2/\delta)}{2n}},
\]
and the corresponding sample size for radius $\epsilon$ is
\[
  n \ge \frac{\log(2/\delta)}{2\epsilon^2}.
\]

For $k$ simultaneously reported Bernoulli rates, a union bound replaces
$\delta$ by $\delta/k$, yielding
\[
  \epsilon(n,k,\delta) =
  \sqrt{\frac{\log(2k/\delta)}{2n}}.
\]
For $m$ guardrails, the number of singleton plus pairwise rates is
$m + m(m-1)/2$.

Computational scaling has a different source. Marginal-only AND and OR bounds
have closed forms. The general finite-atom LP enumerates all $2^m$ atoms, so it
is appropriate for small guardrail stacks, theorem validation, and auditable
examples. Scaling beyond that requires structure, decomposition, or a different
optimization layer.
